% An input fragment for the research paper; no preamble is required.
\subsection{Absorbing the denominator guard in the anchor witness}
\label{subsec:anchor-guard-854}

The explicit construction with $855$ variables has one coordinate whose
only purpose is to invert a product of denominator guards. Its semilocal
anchor has three private witnesses. The first can always be chosen
nonzero, which allows it to absorb the inverse coordinate. This gives an
explicitly specified polynomial in $854$ variables.

\begin{lemma}\label{lem:nonzero-anchor-witness}
Let $C,r\in\mathbb Q$, with $C\ne0$. If
\[
 u^2+6v^2-2w^2-3r^2+2C^4=0
\]
has a rational solution, it has one with $u\ne0$.
\end{lemma}
\begin{proof}
The three witnesses cannot all vanish: that would give
$(r/C^2)^2=2/3$, which is not a rational square. Suppose $u=0$.
If $w\ne0$, replace $(u,w)$ by $(4w,3w)$, leaving $v$ unchanged.
The value of $u^2-2w^2$ is still $-2w^2$, and the new $u$ is nonzero.
If $w=0$, then $v\ne0$. Replace $(u,v)$ by
$(-12v/7,-5v/7)$, leaving $w$ unchanged. Since
\[
 (-12v/7)^2+6(-5v/7)^2=6v^2,
\]
the anchor is again preserved and the new first witness is nonzero.
\end{proof}

\begin{proposition}\label{prop:private-quadratic-guard}
Let $\mathbf x$ be a tuple of variables, and let
$N,G,F_1,\ldots,F_r\in\mathbb Z[\mathbf x]$. Put
\[
 \Psi(\mathbf x,U)=1+N(\mathbf x)G(\mathbf x)^2U^2,
 \qquad
 Q(\mathbf x,U)=\sum_{i=1}^r F_i(\mathbf x)^2+\Psi(\mathbf x,U)^2.
\]
For every $\mathbf a\in\mathbb Q^{\lvert\mathbf x\rvert}$, the following
conditions are equivalent:
\begin{enumerate}
\item $G(\mathbf a)\ne0$, all $F_i(\mathbf a)=0$, and there is a
      nonzero $t\in\mathbb Q$ with $t^2+N(\mathbf a)=0$;
\item there is $U\in\mathbb Q$ with $Q(\mathbf a,U)=0$.
\end{enumerate}
This equivalence also holds over any formally real field.
\end{proposition}
\begin{proof}
If $GU=0$, then $\Psi=1$, so $Q$ cannot vanish over a formally real
field. If $GU\ne0$, the exact identity
\begin{equation}\label{eq:anchor-reciprocal-nonzero}
 \Psi(\mathbf x,U)=(GU)^2\bigl((1/(GU))^2+N\bigr)
\end{equation}
shows that $\Psi=0$ is equivalent to the required quadratic equation
with $t=1/(GU)\ne0$. Conversely, a nonzero root $t$ gives
$U=1/(Gt)$. A sum of squares in a formally real field vanishes
precisely when each summand vanishes, completing the proof.
\end{proof}

\paragraph{Application to the explicit rational system.}
The anchor in the $855$-variable construction is
\[
 u^2+6v^2-2w^2-3r^2+2C^4=0,
 \qquad C=3^{2^{24}},\qquad
 r=f_0^2+f_0f_1-f_1^2.
\]
Here $f_0,f_1$ are rational expressions in point coordinates. None of
$u,v,w$ occurs in any other residual or denominator guard.
Write its cleared numerator as $Bu^2+N$. In the emitted circuit
$B=S^2$ is a polynomial square (indeed, a fourth power), and
$r=R/S$ for a polynomial $R$. The complete cleared identity is
\[
 Bu^2+N=B(u^2+6v^2-2w^2+2C^4)-3R^2.
\]
Let $G$ be the branch's product of all nonconstant denominator factors,
including every factor of $B$. On $G\ne0$, the rescaling $t=Su$
changes the private anchor into $t^2+N=0$. By
Lemma~\ref{lem:nonzero-anchor-witness}, after modifying only the three
private anchor witnesses, one may require $u\ne0$, and hence $t\ne0$.
Let $F_1,\ldots,F_{296}$ be the cleared numerators of all the other
rational equations. Proposition~\ref{prop:private-quadratic-guard}
replaces the anchor and the independent inverse equation by
\[
 \sum_{i=1}^{296}F_i^2+(1+NG^2U^2)^2=0.
\]
Every denominator is nonzero when this equation vanishes. Conversely,
its root recovers $u=1/(GSU)$, with the same values of $v,w$.
The equivalence preserves the projection onto all coordinates outside
$\{u,v,w\}$, and in particular the input coordinate. It need not
preserve each fiber obtained by fixing both $v$ and $w$: the forward
construction may change those two private witnesses.

The formal chart parameter $j$ still ranges over $\{0,\ldots,211\}$,
and the height parameter still ranges over the $256$ signed squarefree
products of $2,3,5,7,11,13,17$. For each of these $54272$ choices form
the polynomial $Q$ above, with that branch's own $N,G,F_i$, and take
the product. A finite product over a field vanishes precisely when
one factor does, so this is the original finite disjunction.
The guards of a branch are required only when that branch supplies
the zero factor. The resulting polynomial has integer coefficients
and uses $854$ total variables, including the input; thus it has
$853$ existential witnesses.

\paragraph{Exact certificate and scope.}
The file \texttt{polynomial-854/polynomial-circuit.json} specifies the
polynomial by an integer arithmetic circuit and the two finite products.
Its $95538$ stored nodes include certificate expressions; $94393$ nodes
are reached from the branch output. Circuit gates are polynomial
expressions, not additional quantified variables. An independent
reader, \texttt{verify\_compact\_854.py}, imports neither producer and
checks the following exact statements: all $296$ untouched residuals
and $1681$ guards agree with the pinned parent circuit; the complete
ternary anchor has the displayed form; $B=S^2$; all three anchor
witnesses are private; every denominator is guarded; the reciprocal
polynomial is $1+NG^2U^2$; the output is the stated sum of $297$
squares; and all declared coordinates occur in the output directed
acyclic graph. Polynomial identities are checked by sparse formal
polynomial arithmetic, not evaluations at sample points.

The independently verified circuit bounds are
\[
 \deg P\le63570781474816,
 \qquad H(P)\le2^{757866260454400},
\]
where $H(P)$ is the largest absolute coefficient of the expanded
polynomial. These are upper bounds from degree and coefficient
$\ell^1$-norm propagation; the expanded polynomial has not been
materialized, and these numbers are not asserted to be its exact
degree or height. The transformation preserves the input projection
of the $855$-variable system. Its identification with $\mathbb Z$
inherits the arithmetic theorem used by that system; the transformation
audit is not an independent verification of every arithmetic argument
in the source manuscript.
